% Ausblick Einheit 4 – Satz vom Nullprodukt – Hauptnummern 47 bis 55
\einheitenkopf[e4][Ausblick 4 Nullprodukt]{Ausblick Einheit 4 von 4 · Satz vom Nullprodukt}
\zweigzeile{Hier lernst du, Gleichungen in Produktform mit dem Satz vom Nullprodukt zu lösen · kommt nächstes Jahr (am Gymnasium schon in Klasse 8 oder 9) · P10 · baut auf: Ausklammern, lineare Gleichungen, p-q-Formel (Einheit 2)}
\setcounter{aufgabe}{46}

\begin{aufgabe}{Ich erkenne, ob rechts eine Null steht.}
Kreuze an, ob der Satz vom Nullprodukt gilt, weil rechts null steht. Löse nichts.
\begin{teile}
\teil $(x - 3) \cdot (x + 1) = 0$ \hfill \kreuz{gilt} \kreuz{gilt nicht}
\teil $x \cdot (x + 2) = 8$ \hfill \kreuz{gilt} \kreuz{gilt nicht}
\teil $(x + 4) \cdot (x - 4) = 0$ \hfill \kreuz{gilt} \kreuz{gilt nicht}
\teil $x \cdot (x - 5) = -4$ \hfill \kreuz{gilt} \kreuz{gilt nicht}
\teil $2x \cdot (x - 1) = 0$ \hfill \kreuz{gilt} \kreuz{gilt nicht}
\end{teile}
\end{aufgabe}

\verfahren{Produkt gleich null}
\begin{aufgabe}{Ich kann eine Gleichung in Produktform mit dem Satz vom Nullprodukt lösen.}
Löse die Gleichung.
\rechenplatz{2}
\begin{gleichungsraster}[2]
\gl{(x - 2) \cdot (x - 5) = 0} & \gl{(x - 1) \cdot (x - 6) = 0} \\
\gl{(x - 3) \cdot (x - 4) = 0} & \gl{(x - 7) \cdot (x - 2) = 0} \\
\gl{(x + 5) \cdot (x - 1) = 0} & \gl{x \cdot (x + 6) = 0} \\
\gl{(x - 8) \cdot (x - 8) = 0} & \\
\end{gleichungsraster}
\end{aufgabe}

\begin{aufgabe}{Ich kann prüfen, welche Zahl eine Gleichung in Produktform erfüllt.}
Setze die Zahl ein und kreuze an: wahre Aussage (wA) oder falsche Aussage (fA).
\begin{teile}
\teil $x = 3$ \quad in \quad $(x - 3) \cdot (x + 2) = 0$ \hfill \kreuz{wA} \kreuz{fA}
\teil $x = -2$ \quad in \quad $x \cdot (x + 4) = -4$ \hfill \kreuz{wA} \kreuz{fA}
\teil $x = 2$ \quad in \quad $x \cdot (x + 4) = -4$ \hfill \kreuz{wA} \kreuz{fA}
\teil $x = -1$ \quad in \quad $(x + 3) \cdot (x - 2) = -6$ \hfill \kreuz{wA} \kreuz{fA}
\teil $x = 1$ \quad in \quad $(x - 4) \cdot (x + 2) = 9$ \hfill \kreuz{wA} \kreuz{fA}
\anweisung{Kreuze den Wert an, der die Gleichung erfüllt. (P10 2025 OS)}
\teil $x \cdot (x + 7) = -10$ \hfill \kreuz{$x = 2$} \kreuz{$x = 5$} \kreuz{$x = -2$} \kreuz{$x = -10$}
\end{teile}
\end{aufgabe}

\verfahren{Ausklammern}
\begin{aufgabe}{Ich kann eine Gleichung ohne Zahlglied durch Ausklammern lösen.}
Löse die Gleichung.
\rechenplatz{3}
\begin{gleichungsraster}[3]
\gl{x^2 + 4x = 0} & \gl{x^2 + 9x = 0} \\
\gl{x^2 - x = 0} & \gl{3x^2 - 12x = 0} \\
\gl{x^2 = 5x} & \\
\end{gleichungsraster}
\end{aufgabe}

\verfahren{Produkt ungleich null}
\begin{aufgabe}{Ich kann eine Produktgleichung lösen, bei der rechts nicht null steht.}
Löse die Gleichung.
\rechenplatz{4}
\begin{gleichungsraster}[3]
\gl{x \cdot (x + 2) = 35} & \gl{(x - 2) \cdot (x + 3) = 14} \\
\gl{x \cdot (x - 3) = 10} & \\
\end{gleichungsraster}
\end{aufgabe}

\begin{aufgabe}{Ich kann eine Gleichung lösen, indem ich die linke Seite als Produkt schreibe.}
Schreibe die linke Seite als Produkt zweier Klammern und löse.
\begin{gleichungsraster}[2]
\gl{x^2 + 8x + 15 = 0} & \gl{x^2 - 7x + 12 = 0} \\
\gl{x^2 + 2x - 3 = 0} & \\
\anweisung{Berechne die Nullstellen der Parabel durch Faktorisieren. (P10 2020 OS)}
\gl{y = 2x^2 + 14x + 20} & \\
\end{gleichungsraster}
\end{aufgabe}

\begin{aufgabe}{Ich finde den Fehler beim Satz vom Nullprodukt.}
Finde den Fehler, benenne ihn und korrigiere die Rechnung.
\begin{teile}
\teil Nora rechnet: \rechnung{x^2 &= 3x &&\mid : x \\ x &= 3}
\teil Ali rechnet: \rechnung{(x + 2) \cdot (x - 6) &= 0 \\ x_1 &= 2, \quad x_2 = -6}
\teil Kim rechnet: \rechnung{x \cdot (x + 1) &= 12 \\ x = 12 \quad &\text{oder} \quad x + 1 = 12 \\ x_1 = 12 \quad &\text{oder} \quad x_2 = 11}
\end{teile}
\end{aufgabe}

\begin{aufgabe}{Ich kann begründen, wann der Satz vom Nullprodukt gilt.}
Begründe.
\begin{teile}
\teil Warum ist $(x - 4) \cdot (x + 1)$ gleich null, wenn $x = 4$ ist?
\teil Warum darf man $x^2 = 6x$ nicht durch $x$ teilen?
\end{teile}
\end{aufgabe}

\begin{aufgabe}{Ich kann mit dem Satz vom Nullprodukt eine Sachfrage beantworten.}
Ein Ball wird vom Boden senkrecht nach oben geworfen. Nach $t$ Sekunden ist er $h = 15t - 5t^2$ Meter hoch.
\begin{teile}
\teil Schreibe $15t - 5t^2$ als Produkt.
\teil Nach wie vielen Sekunden ist der Ball wieder am Boden?
\teil Was bedeutet die Lösung $t = 0$ im Sachzusammenhang?
\end{teile}
\end{aufgabe}
